% Results section — DESI paper v3 (complete rewrite)
% Narrative: validation → AFR>EAS → bottleneck scan → robustness → AMR

\section*{Results}

\subsection*{DESI accurately recovers pairwise coalescent depth across demographic models}

We validated the Pairwise Windowed Likelihood (PWL) method on coalescent simulations spanning three demographic scenarios: a standard Out-of-Africa model (Model OoA: ancestral $N_e = 15{,}000$, African $N_e = 15{,}000$, East Asian $N_e = 3{,}000$ post-bottleneck at $\sim$65 ka~\citep{Gutenkunst2009}); a strictly panmictic single-population model (Model A: constant $N_e = 10{,}000$, group labels randomly assigned); and a reduced-ancestral-$N_e$ model (Model D: same OoA structure but ancestral $N_e = 12{,}000$). Across all three models and three independent replicates (seeds 42/123/456), PWL recovered mean within-group TMRCA depths within 3.5\% of the exact tree-sequence values computed from \texttt{msprime}~\citep{Kelleher2016} (Fig.~\ref{fig:model_compare}; Table~\ref{tab:calibration}). No systematic bias towards inflating between-group differences was observed: under Model A (randomly labelled groups), the PWL-inferred AFR$-$EAS TMRCA difference was $0.5 \pm 0.3$ ka, confirming the method does not artifactually manufacture structure when none exists. Under Model D, the AFR$-$EAS difference was 396.5 ka against an exact value of 370.1 ka ($+7.1\%$), a modest overestimate that is conservative relative to our main claim. Having established calibration accuracy, we applied DESI to the 1000 Genomes Project.

\subsection*{Within-African populations show $\sim$1.4$\times$ deeper genealogical ancestry than East Asians genome-wide}

Across 27,507 valid 100 kb windows spanning all 22 autosomes, within-AFR mean pairwise TMRCA was $1{,}229.5 \pm 6.4$ ka (block-jackknife SE), substantially deeper than within-EAS ($858.8 \pm 5.4$ ka), within-EUR ($920.2 \pm 5.5$ ka), within-SAS ($984.8 \pm 5.6$ ka), and within-AMR ($950.2 \pm 5.6$ ka) (Fig.~\ref{fig:tmrca_matrix}a; Table~\ref{tab:tmrca_means}). Between-group comparisons involving AFR are similarly deep (AFR$\times$EAS: $1{,}237.7 \pm 6.6$ ka), while all non-African between-group comparisons cluster within a narrow $962$--$995$ ka range, indicating a shared recent ancestral genealogy for the non-African clade.

The within-AFR minus within-EAS difference is $370.7 \pm 8.4$ ka ($Z = 44$). This difference is highly consistent across chromosomes: all 22 autosomes show AFR $>$ EAS ($Z = 73$ from a leave-one-chromosome-out jackknife; Fig.~\ref{fig:tmrca_matrix}b), with inter-chromosomal coefficient of variation of only 6.4\% (range: 340.9--440.0 ka). The signal is spatially uniform across the genome rather than concentrated in particular chromosomal regions.

\subsection*{FitCoal's bottleneck severity is quantitatively inconsistent with the observed TMRCA distribution}

To directly test whether FitCoal's inferred bottleneck parameterization~\citep{Hu2023} is consistent with the observed within-AFR TMRCA distribution, we designed a parameter scan. We simulated 112 demographic scenarios spanning a grid of bottleneck severity ($N_e^{\text{bn}} \in \{300, 500, 800, 1{,}280, 2{,}000, 4{,}000, 8{,}000, 20{,}000\}$) and duration ($d \in \{30, 50, 80, 117, 150, 200, 300\}$ ka), with the bottleneck epoch fixed at 813--$(813 + d)$ ka following FitCoal's timing, under two ancestral $N_e$ values (20,000 and 50,000). All scenarios share the same OoA structure as the real data (AFR modern $N_e = 15{,}000$; EAS $N_e = 3{,}000$ post-bottleneck at 65 ka). For each scenario, we applied the identical per-window TMRCA analysis using tskit branch-mode diversity (equivalent to DESI's group-mean TMRCA), and computed two summary statistics matched against the empirical observations:

\textbf{(i) $P(\bar{T}_\text{AFR} > 930\,\text{ka})$}: the fraction of 100 kb windows for which the within-AFR group-mean TMRCA exceeds 930 ka. This statistic is sensitive to whether genealogical depth is concentrated near the bottleneck epoch (predicting lower $P$) or distributed more broadly (higher $P$).

\textbf{(ii) Mean within-AFR TMRCA}: the genome-wide mean $\bar{T}_\text{AFR}$.

\paragraph{FitCoal exact parameterization is incompatible with observations.}
Using FitCoal's own ancestral $N_e = 50{,}000$, the exact FitCoal parameters ($N_e^{\text{bn}} = 1{,}280$, $d = 117$ ka) predict $P(\bar{T}_\text{AFR} > 930\,\text{ka}) = 0.615$ (from 600 simulated windows), compared to the observed value of 0.707 ($Z = 4.6$, $p < 0.0001$; Fig.~\ref{fig:scan}c). The FitCoal scenario also predicts mean within-AFR TMRCA of 1,105 ka, underestimating the observed 1,229 ka by 10.2\% (Fig.~\ref{fig:scan}e, red star). Both discrepancies are in the same direction: FitCoal predicts genealogies that are systematically younger and less likely to exceed 930 ka than the data show. 

Mechanistically, the FitCoal bottleneck ($N_e = 1{,}280$ for 117 ka) has a theoretical per-pair coalescence probability of $1 - e^{-4179/2560} \approx 80.5\%$ within the 813--930 ka window, forcing the majority of genealogies into the bottleneck epoch and depressing the fraction of windows above 930 ka. The observed $P = 0.707$ is incompatible with this prediction at $Z = 4.6$.

\paragraph{The compatible parameter space requires a much weaker bottleneck.}
The parameter scan identifies which combinations of $N_e^{\text{bn}}$ and duration simultaneously satisfy $|P_\text{sim} - P_\text{obs}| < 0.10$ and $|\bar{T}_\text{sim} - \bar{T}_\text{obs}| < 200$ ka. The closest match in the ancestral $N_e = 50{,}000$ space is $N_e^{\text{bn}} \approx 2{,}000$ with $d \approx 150$ ka (predicted $P = 0.778$, $\bar{T}_\text{AFR} = 1{,}239$ ka; Fig.~\ref{fig:scan}a,b,f). This represents a bottleneck approximately 1.6$\times$ less severe than FitCoal's reported value. The full heatmaps (Fig.~\ref{fig:scan}a,b) show a clear gradient: stronger or longer bottlenecks (lower $N_e^{\text{bn}}$ or larger $d$) progressively shift the simulation statistics away from the observations, with FitCoal's exact point (red box) lying outside the compatible region. Notably, neither the ``no bottleneck'' scenario ($N_e^{\text{bn}} = N_e^{\text{anc}}$) nor FitCoal's extreme parameterization ($N_e^{\text{bn}} = 1{,}280$) provides a compatible fit; the data favor an intermediate severity.

We emphasize that this analysis tests FitCoal's specific panmictic parameterization against empirical TMRCA summary statistics. It does not rule out a population reduction event at $\sim$930 ka, nor does it adjudicate between a weaker panmictic bottleneck and partial ancestral population structure, as both alternatives could produce the observed intermediate $P$ values.

\subsection*{The genealogical depth asymmetry is robust to background selection, mutation rate, and genomic architecture}

\paragraph{Background selection.}
To assess whether differential background selection (BGS) between AFR and EAS populations could drive the observed TMRCA asymmetry, we restricted analysis to 30.7\% of autosomal windows located $>$50 kb from any annotated gene. In these putatively neutral windows, within-AFR TMRCA was $1{,}197.6 \pm 14.9$ ka and within-EAS was $836.6 \pm 11.7$ ka, a difference of $361.1 \pm 19.0$ ka ($Z = 19.0$; Fig.~\ref{fig:robustness}a) --- only 2.6\% smaller than the genome-wide estimate (370.7 ka). Additionally, chromosome 19 (the most gene-dense autosome) shows the \textit{largest} per-chromosome AFR$-$EAS difference (440.0 ka) rather than the smallest expected under the BGS hypothesis. These two converging lines of evidence exclude BGS as a meaningful contributor.

\paragraph{Mutation rate invariance.}
The AFR/EAS TMRCA ratio is $1{,}229.5/858.8 = 1.432$. Because all TMRCA estimates scale as $\hat{T} \propto 1/\mu$, this ratio is analytically invariant to the assumed mutation rate; we confirm this numerically across $\mu \in \{1.0, 1.2, 1.4, 1.6\} \times 10^{-8}$ per bp per generation (Fig.~\ref{fig:robustness}b). All model-comparison conclusions are therefore independent of mutation-rate calibration.

\paragraph{Chromosomal architecture.}
Acrocentric and metacentric chromosomes show statistically indistinguishable AFR$-$EAS differences (370.9 vs.\ 373.4 ka; difference: 2.5 ka). Chromosome length is not a significant predictor of the AFR$-$EAS difference (Spearman $\rho = -0.37$, $p = 0.09$). No chromosomal confound accounts for the pattern.

\subsection*{AMR sub-population stratification validates genealogical depth tracking}

Admixed American populations provide a positive-control test: populations with predominantly African ancestry should show within-group TMRCA similar to AFR, while those with predominantly indigenous American (East-Asian-like) ancestry should match EAS. Within-ACB (African-Caribbean; predominantly African ancestry) was $1{,}252.8 \pm 11.3$ ka, indistinguishable from within-AFR ($\Delta = 23$ ka, $<2 \times$ SE; Fig.~\ref{fig:robustness}c). Within-PEL (Peruvian; predominantly indigenous Andean ancestry) was $886.7 \pm 70.6$ ka, closely matching within-EAS. Within-PUR (Puerto Rican; tripartite admixture) was intermediate at $1{,}029.0 \pm 52.9$ ka. This gradient from AFR-like to EAS-like TMRCA, tracking known ancestry proportions, confirms that DESI estimates reflect genuine deep genealogical signal rather than methodological artifacts from recent population structure.
